
\documentclass{article}
\usepackage{colortbl}
\usepackage{makecell}
\usepackage{multirow}
\usepackage{supertabular}

\begin{document}

\newcounter{utterance}

\centering \large Interaction Transcript for game `cladder', experiment `full\_v1.5\_default', episode 2504 with Qwen3.5{-}4B.
\vspace{24pt}

{ \footnotesize  \setcounter{utterance}{1}
\setlength{\tabcolsep}{0pt}
\begin{supertabular}{c@{$\;$}|p{.15\linewidth}@{}p{.15\linewidth}p{.15\linewidth}p{.15\linewidth}p{.15\linewidth}p{.15\linewidth}}
    \# & \multicolumn{2}{c}{Player} && \multicolumn{2}{c}{Game Master} \\
    \hline

    \theutterance \stepcounter{utterance}  
    & & & \multicolumn{4}{p{0.6\linewidth}}{
        \cellcolor[rgb]{0.9,0.9,0.9}{
            \makecell[{{p{\linewidth}}}]{
                \texttt{\tiny{[P1$\langle$GM]}}
                \texttt{You are an expert in causal inference. The following question is not a typical commonsense query, but rather a meticulously designed question created by a professor specializing in causal inference, intended to assess the students' mastery of the course content.} \\
\texttt{Imagine a self{-}contained, hypothetical world with only the following conditions, and without any unmentioned factors or causal relationships: Kidney stone size has a direct effect on treatment and recovery. Treatment has a direct effect on recovery. For patients who have small kidney stones and not receiving treatment, the probability of recovery is 88\%. For patients who have small kidney stones and receiving treatment, the probability of recovery is 72\%. For patients who have large kidney stones and not receiving treatment, the probability of recovery is 16\%. For patients who have large kidney stones and receiving treatment, the probability of recovery is 4\%. The overall probability of large kidney stone is 41\%. For patients receiving treatment, would it be less likely to see recovery if the patient had received no treatment?} \\
\texttt{Start your answer with 'Yes' or 'No', followed by additional reasoning or evidence to support your explanation.} \\
            }
        }
    }
     \\ \\

    \theutterance \stepcounter{utterance}  
    & \multicolumn{4}{p{0.6\linewidth}}{
        \cellcolor[rgb]{0.561,0.737,0.561}{
            \makecell[{{p{\linewidth}}}]{
                \texttt{\tiny{[P1$\rangle$GM]}}
                \texttt{Yes} \\
\\ 
\texttt{**Reasoning:**} \\
\\ 
\texttt{To answer this question, we must distinguish between the **causal effect** of treatment (what would happen if we intervened) and the **observational association** (what we see in the data). The question asks: "For patients receiving treatment, would it be less likely to see recovery if the patient had received no treatment?"} \\
\\ 
\texttt{This phrasing asks us to compare the counterfactual outcome for a specific group: specifically, the probability of recovery for patients who *actually received treatment* but had *not received treatment* (denoted as \$P(\textbackslash{}text\{Recovery\} | \textbackslash{}text\{Treatment\}, \textbackslash{}text\{No Treatment\}\$).} \\
\\ 
\texttt{1.  **Identify the Causal Effect**: The problem states that "Treatment has a direct effect on recovery." We are given the conditional probabilities of recovery based on stone size and treatment status:} \\
\texttt{    *   Small stones + Treatment: 72\%} \\
\texttt{    *   Small stones + No Treatment: 88\%} \\
\texttt{    *   Large stones + Treatment: 4\%} \\
\texttt{    *   Large stones + No Treatment: 16\%} \\
\\ 
\texttt{2.  **Analyze the Causal Impact**:} \\
\texttt{    *   For patients with **small** stones, moving from "No Treatment" to "Treatment" changes the recovery probability from 88\% to 72\%. This is a **decrease** in recovery probability. Therefore, for this subgroup, receiving treatment makes recovery *less* likely than not receiving it.} \\
\texttt{    *   For patients with **large** stones, moving from "No Treatment" to "Treatment" changes the recovery probability from 16\% to 4\%. This is also a **decrease** in recovery probability.} \\
\\ 
\texttt{3.  **Evaluate the Counterfactual for the "Treatment" Group**:} \\
\texttt{    The question asks about the hypothetical scenario where patients who *received treatment* had instead received *no treatment*.} \\
\texttt{    *   If a patient with small stones received treatment, their recovery probability is 72\%. If they had received no treatment, it would be 88\%. Thus, removing treatment increases recovery probability.} \\
\texttt{    *   If a patient with large stones received treatment, their recovery probability is 4\%. If they had received no treatment, it would be 16\%. Thus, removing treatment increases recovery probability.} \\
\\ 
\texttt{    In both subgroups, the causal effect of treatment is negative (it lowers the probability of recovery). Consequently, for any patient who received treatment, the counterfactual outcome where they received no treatment would result in a **higher** probability of recovery (88\% > 72\% and 16\% > 4\%).} \\
\\ 
\texttt{4.  **Address the Specific Wording**:} \\
\texttt{    The question asks: "would it be **less likely** to see recovery if the patient had received no treatment?"} \\
\texttt{    *   Current state (with treatment): 72\%} \\
            }
        }
    }
    & & \\ \\

    \theutterance \stepcounter{utterance}  
    & & & \multicolumn{2}{p{0.3\linewidth}}{
        \cellcolor[rgb]{0.9,0.9,0.9}{
            \makecell[{{p{\linewidth}}}]{
                \texttt{\tiny{[GM$|$GM]}}
                \texttt{Yes} \\
            }
        }
    }
    & & \\ \\

    \theutterance \stepcounter{utterance}  
    & & & \multicolumn{2}{p{0.3\linewidth}}{
        \cellcolor[rgb]{0.9,0.9,0.9}{
            \makecell[{{p{\linewidth}}}]{
                \texttt{\tiny{[GM$|$GM]}}
                \texttt{no} \\
            }
        }
    }
    & & \\ \\

    \theutterance \stepcounter{utterance}  
    & & & \multicolumn{2}{p{0.3\linewidth}}{
        \cellcolor[rgb]{0.9,0.9,0.9}{
            \makecell[{{p{\linewidth}}}]{
                \texttt{\tiny{[GM$|$GM]}}
                \texttt{game\_result = LOSE} \\
            }
        }
    }
    & & \\ \\

\end{supertabular}
}

\end{document}
